% ECONOMICS_PROBLEM: {"id": "J65-96", "title": "Unemployment-insurance design", "jel_code": "J65", "source_id": "OP-0478"}
% FIRST_POSTED: 2026-10-11
% ATTEMPT_STATUS: solved
% ENTRY_KIND: research_attempt
\documentclass[11pt]{article}
\usepackage[a4paper,margin=24mm]{geometry}
\usepackage[T1]{fontenc}
\usepackage{lmodern,microtype,amsmath,amssymb,amsthm,mathtools}
\usepackage{enumitem,xurl,hyperref}
\hypersetup{colorlinks=true,linkcolor=blue,urlcolor=blue,citecolor=blue}
\setlength{\parskip}{3pt}
\newtheorem{theorem}{Theorem}
\newtheorem{lemma}{Lemma}
\newtheorem{proposition}{Proposition}
\newcommand{\E}{\mathbb E}
\newcommand{\R}{\mathbb R}
\newcommand{\ind}{\mathbf1}
\title{Welfare derivatives for equilibrium unemployment insurance}
\author{JiHo Bae}
\date{October 11, 2026}
\begin{document}
\maketitle

\paragraph{Authorship and provenance.}
Author: JiHo Bae. Prepared with GPT Codex assistance; the precise serving
revision was not exposed. The complete original argument is retained in
the \href{https://github.com/jbaelaw/equilibrium-unemployment-insurance-proof/blob/3de5842fb6f6686849363cd4c9f8dc32a395d02c/paper/Proof.tex}{immutable TeX source} and
\href{https://github.com/jbaelaw/equilibrium-unemployment-insurance-proof/blob/3de5842fb6f6686849363cd4c9f8dc32a395d02c/paper/Proof.pdf}{proof PDF}. This submission adds catalogue metadata,
provenance and a completion statement to that argument. Its preparation
included exact-target and primary-source checks and separate mathematical
reviews. This is a complete proof claim submitted for independent expert
review; community acceptance remains pending.

\begin{abstract}
We characterize local benefit reforms in a declared class of infinite-horizon labor-search economies. Benefits are financed each period, workers and firms optimize globally, and the initial population distribution is fixed. At every interior regular equilibrium, one finite linear system gives the welfare gradient, including changes in search, separations, vacancies, dividends and taxes. When welfare is evaluated at the worst equilibrium, its directional derivative is the minimum of these gradients over the baseline worst equilibria. An explicit economy with two active workers and a hiring firm proves that the assumptions hold on a nonempty open set of primitives.
\end{abstract}

\section{The target and the economy}
Catalogue problem J65-96 expressly permits a characterization of local reform welfare derivatives in a declared primitive class. We give that characterization below. The source review motivates duration-sensitive benefits and equilibrium effects \cite[\S\S3.6.2, 3.7.1, 3.9]{ui}; it does not supply our model or theorem.

There are finitely many workers $i\in I$ and firms $f\in K$. The public Markov state $x$ belongs to a finite set $X$ and includes employment, employer, unemployment duration, observable worker characteristics and aggregate conditions. Duration is the actual elapsed spell length. The primitives impose a finite maximum spell length $D$, for example by guaranteed reemployment at $D$; no tail of longer spells is silently collapsed. Liquid assets and borrowing are fixed at zero. Firm ownership shares are fixed endowments, and there is no asset trading. These are restrictions defining the class, rather than approximations to an unrestricted asset economy.

A policy $b$ is a vector in a bounded open benefit neighborhood $B\subset\R^q$. Its coordinates specify duration- and characteristic-dependent benefits. Worker $i$'s receipt in state $x$ is a fixed linear coordinate selection $b_i(x)$, equal to zero when employed. Fixed tax shares $\gamma_i\ge0$ sum to one, and
\[
 t_i(x,b)=\gamma_i\sum_{j\in I}b_j(x).
\]
Thus taxes finance benefits in every state. The initial law $\mu_0$ on $X$, discount $0<\beta<1$, and nonnegative welfare weights $\omega_i$ with $\sum_i\omega_i=1$ are fixed across reforms.

Each worker chooses search or employment-retention effort $a_i(x)$ in a compact box; each firm chooses vacancy effort $a_f(x)$ in a compact box. Some states may have no choice. All boxes are independent of $b$. The specified transition kernel $P(x'\mid x,a)$ incorporates matching, offers and separations, respects employment and duration accounting, and is continuously differentiable twice in actions. For each player separately it is affine in that player's entire action vector, with the other actions fixed. It is stochastic throughout the action boxes. All relevant information is contained in $x$.

Output $y_f(x)$, wage bills $w_f(x)$, individual wages $w_i(x)$ and outside endowments $e_i(x)$ are primitives, with $\sum_i w_i(x)=\sum_f w_f(x)$. Firm $f$ earns
\[
 \pi_f(x,a_f)=y_f(x)-w_f(x)-k_f(x,a_f).
\]
The twice continuously differentiable effort costs $k_f$ and $k_i$ have uniformly positive definite Hessians. Ownership shares $\alpha_{if}\ge0$ satisfy $\sum_i\alpha_{if}=1$. Consumption and period payoffs are
\begin{align}
 c_i(x,a,b)&=e_i(x)+w_i(x)+b_i(x)-t_i(x,b)
                   +\sum_f\alpha_{if}\pi_f(x,a_f),\label{eq:consumption}\\
 r_i(x,a,b)&=u_i(c_i(x,a,b))-k_i(x,a_i),\qquad
 r_f(x,a,b)=\pi_f(x,a_f).\label{eq:flows}
\end{align}
Here $u_i$ is twice continuously differentiable and increasing. Consumption is uniformly positive for all action profiles and every $b\in B$. Payoffs and their derivatives are bounded on this domain. Firms maximize discounted profit; workers maximize discounted utility. Feasible consumption is the amount in \eqref{eq:consumption}, since there is no saving. The accounts give, state by state,
\[
 \sum_i c_i=\sum_i e_i+\sum_f y_f-\sum_f k_f.
\]
Wages, benefits and taxes cancel, and distributed profits appear once. The transition technology clears jobs; the displayed identity clears goods and the tax rule clears the government account.

We take equilibrium to mean stationary Markov-perfect Nash equilibrium. A unilateral deviation may nevertheless be randomized and depend on the entire observed history. Let $\mathcal E(b)$ contain \emph{all} equilibria under this convention. For equilibrium $e$, write $V_i^e(x)$ for discounted worker utility, and define
\begin{equation}\label{eq:objective}
 J(b)=\min_{e\in\mathcal E(b)}W(e,b),\qquad
 W(e,b)=\sum_i\omega_i\sum_x\mu_0(x)V_i^e(x).
\end{equation}
The catalogue's policy value in this class is $\sup_{b\in B}J(b)$. Our target is the local reform derivative of $J$, rather than a universal globally optimal benefit rule.

\section{Global optimization and the equilibrium equations}
\begin{lemma}\label{lem:equilibrium}
The equilibrium set is nonempty and compact. Every stationary behavioral equilibrium is pure. With opponents' stationary policies fixed, each player's stationary best response is unique and is optimal from every state against all history-dependent deviations.
\end{lemma}
\begin{proof}
For a player $p$, fix its opponents' policies. Its Bellman operator sends $V$ to
\[
 (TV)(x)=\max_{a_p}\left\{r_p(x,a_p,a_{-p}(x),b)
        +\beta\sum_{x'}P(x'\mid x,a_p,a_{-p}(x))V(x')\right\}.
\]
Stochasticity makes this a contraction of modulus $\beta$ in the supremum norm, so it has a unique fixed point. The expression maximized is strictly concave in $a_p$: continuation is affine, and current payoff has strictly negative definite own-action Hessian. For workers their consumption does not depend on their own effort. Thus the maximizing action in every state is unique.

For randomized stationary opponents, take expectations over their actions. This preserves uniform own-action strict concavity and affine continuation, so the same unique-maximizer argument applies. Nonlinear utility is evaluated before taking that expectation.

The Bellman inequality, iterated along any randomized history-dependent deviation, bounds its expected payoff through date $T$ plus the discounted terminal value by the initial Bellman value. Boundedness makes the terminal term vanish as $T\to\infty$. The stationary maximizing policy attains equality. This proves global optimality and excludes nondegenerate mixing at any state of a stationary behavioral equilibrium, including equilibria with randomized opponents.

The Bellman fixed point varies continuously with opponents' policies: uniform changes in the operator change its fixed point by at most the operator change divided by $1-\beta$. Compactness and uniqueness of each maximizing action then imply continuity of the best-response policy. Brouwer's theorem \cite{brouwer} applied to the product of the action boxes gives a fixed point. The fixed-point set is closed in that compact product. Values depend continuously on policies, so the set of equilibria including their values is compact.
\end{proof}

Let $z$ collect all players' value vectors $V_p$ and all active action coordinates $a_p(x)$. At an interior equilibrium define the square system $F(z,b)=0$ by, for every player $p$ and state $x$,
\begin{align}
 F^V_{p,x}&=V_p(x)-r_p(x,a(x),b)
              -\beta\sum_{x'}P(x'\mid x,a(x))V_p(x'),\label{eq:bellman}\\
 F^A_{p,x}&=D_{a_p}r_p(x,a(x),b)
              +\beta\sum_{x'}D_{a_p}P(x'\mid x,a(x))V_p(x').\label{eq:foc}
\end{align}
There is one equation in \eqref{eq:foc} for each active action coordinate. Strict concavity proves that these first-order equations give global Bellman maximizers. Consequently every feasible interior root is an equilibrium, and every interior equilibrium is a root. This equivalence includes firm optimization, not just workers' first-order conditions.

\section{The full local derivative}
Fix a baseline $b_0$ at which \emph{every} equilibrium is interior and the matrix $F_z$ is nonsingular at every equilibrium. These are the regularity assumptions on the declared class of baselines. We do not remove boundary or singular equilibria from a larger correspondence. Section~4 proves that the assumptions hold for an open set of labor-market primitives.

For each baseline equilibrium $e$, let $\lambda^e$ solve
\begin{equation}\label{eq:adjoint}
 F_z(z^e,b_0)^\top\lambda^e=\eta,\qquad
 \eta^\top z=\sum_i\omega_i\sum_x\mu_0(x)V_i(x).
\end{equation}
Write $\lambda^e_{i,x}$ for its worker Bellman-equation coordinates, as distinct from its action-equation and firm coordinates. A reform direction $h\in\R^q$ induces receipts $h_i(x)$ through the same coordinate selections as $b_i(x)$.

\begin{theorem}\label{thm:derivative}
There are finitely many baseline equilibria. Put
\[
 E_* =\{e\in\mathcal E(b_0):W(e,b_0)=J(b_0)\}.
\]
The one-sided directional derivative of worst-equilibrium welfare exists for every $h$ and is
\begin{equation}\label{eq:derivative}
 DJ(b_0;h)=\min_{e\in E_*}
 \sum_{i\in I,x\in X}\lambda^e_{i,x}u_i'(c_i^e(x))
       \left[h_i(x)-\gamma_i\sum_jh_j(x)\right].
\end{equation}
All changes in search, retention, vacancies, profits and their distribution are included through \eqref{eq:adjoint}. If the right side is positive, the reform strictly increases $J$ for every sufficiently small positive step. A negative right side strictly decreases it. The formula is a derivative characterization from primitives and equilibrium quantities; it makes no claim that unemployment-duration data alone identify its coefficients.
\end{theorem}
\begin{proof}
Nonsingularity and the implicit-function theorem isolate every baseline equilibrium. A compact set consisting of isolated roots is finite: any infinite subset would have an accumulation point in the equilibrium set, contradicting isolation. Each equilibrium extends to a unique continuously differentiable branch near $b_0$.

These branches cover \emph{all} nearby equilibria. Otherwise choose benefit vectors $b_n\to b_0$ with equilibria outside their branch neighborhoods. Equilibrium action profiles lie in compact boxes, and values are uniformly bounded. A subsequence converges, and continuity of the best-response equations makes its limit a baseline equilibrium. Its interior neighborhood is covered by the corresponding implicit-function branch, a contradiction. Conversely each branch stays feasible and interior, and \eqref{eq:bellman}--\eqref{eq:foc} verify that it remains an equilibrium.

On a branch, differentiation gives $z'[h]=-F_z^{-1}F_bh$. Since welfare is $\eta^\top z$, its derivative is $-\lambda^\top F_bh$. The kernel and effort costs are independent of $b$, while current worker consumption is independent of the worker's own action. Hence the only nonzero coordinates of $F_bh$ are
\[
 (F_bh)^V_{i,x}=-u_i'(c_i^e(x))
            \left[h_i(x)-\gamma_i\sum_jh_j(x)\right].
\]
This gives the branch expression in \eqref{eq:derivative}.

Finally $J$ is the minimum of the finitely many branch welfare functions. Nonminimizing branches have a strictly positive baseline welfare gap and stay inactive locally. Taylor expansion of the remaining finitely many branches, with a common $o(t)$ remainder, gives
\[
 J(b_0+th)=J(b_0)+t\min_{e\in E_*}DW_e(b_0)[h]+o(t)
 \qquad(t\downarrow0).
\]
This proves the formula and its strict-sign implications. The fixed vector $\mu_0$ in \eqref{eq:adjoint} was never replaced by a policy-dependent stationary distribution.
\end{proof}

At a unique equilibrium the derivative is linear in $h$. At tied worst equilibria it may have a corner. One must use the minimum over the \emph{baseline worst} equilibria, rather than a selected branch or all branches indiscriminately.

\section{A nonempty open class with active labor-market responses}
\begin{proposition}\label{prop:example}
The hypotheses of Theorem~\ref{thm:derivative} hold on a nonempty open parameter set with two workers, one firm, endogenous search, endogenous retention and positive vacancy effort.
\end{proposition}
\begin{proof}
Each worker has status $E,U_1,U_2$, giving nine joint states. Status $U_d$ means an actual unemployment spell of $d$ periods. Set $B(E)=E$, $B(U_1)=U_2$, $B(U_2)=E$, componentwise. The last transition is guaranteed reemployment after two periods. Let $H_i(x)$ replace worker $i$'s component of $B(x)$ by $E$, and let $S_i(x)$ replace it by $U_1$.

A worker at $U_1$ chooses search $s_i\in[0,1]$, and a worker at $E$ chooses retention $r_i\in[0,1]$. There is no effort choice at $U_2$. The firm chooses $v\in[0,1]$ exactly in states containing a $U_1$ worker. For each $U_1$ worker assign probability $0.05s_i(1+0.1v)$ to $H_i(x)$; for each employed worker assign probability $0.1(1-r_i)$ to $S_i(x)$. Assign the remaining probability to $B(x)$. These events are exclusive, their total probability is at most $0.31$, and the kernel is affine in each player's own action. Duplicate next states simply have their probabilities added. Employment and duration transitions are exact.

Wages are $w_1=2,w_2=3$ when employed; corresponding output is $w_i+1$. Outside endowments are $L_1=10,L_2=12$. Both tax shares and both ownership shares are $1/2$. Baseline benefits at both unemployment durations are one. Worker utility is $\log c_i$, and active effort cost is $s_i^2/2$ or $r_i^2/2$. Vacancy cost is $v^2/2$ in its active states and zero elsewhere. Writing $E_i=\ind\{x_i=E\}$, firm profit and worker consumption at the baseline are
\[
 \pi=E_1+E_2-v^2/2,\qquad
 c_i=L_i+w_iE_i+E_j-v^2/4\quad(j\ne i),
\]
where $v=0$ in states without a vacancy choice. Consumption is strictly positive, including under nearby benefit and primitive perturbations.

First temporarily put $\beta=0$. Each worker uniquely chooses zero effort and the firm uniquely chooses zero vacancy effort. Thus there is a unique equilibrium. At this point the value-equation Jacobian with respect to values is the identity; its action derivatives vanish. Each action-equation derivative with respect to its own action is $-1$, and all its other derivatives vanish. Therefore $F_z$ is nonsingular.

As $\beta\downarrow0$, every equilibrium converges to this point. Indeed the Bellman tail is uniformly $O(\beta)$; any subsequential limit must maximize each player's current payoff and therefore has all actions zero and their prescribed current values. At zero discount the continuation-value difference for worker $i$ between $H_i(x)$ and $B(x)$, or between $B(x)$ and $S_i(x)$, is
\[
 \log\frac{L_i+w_i+E_j(B(x))}{L_i+E_j(B(x))}>0.
\]
The firm's difference between $H_i(x)$ and $B(x)$ is one. These strict inequalities persist uniformly at all equilibria for sufficiently small positive $\beta$.

The derivative of a worker's Bellman objective at zero search is consequently positive, as is its derivative at zero retention. At effort one the derivative is $-1+O(\beta)<0$. All worker choices are therefore interior. In a state with a $U_1$ worker, the firm's derivative at $v=0$ is
\[
 0.005\beta\sum_{i:x_i=U_1}s_i
              [V_f(H_i(x))-V_f(B(x))]>0.
\]
At $v=1$ it is $-1+O(\beta)<0$, so vacancy effort is also interior and positive. The implicit-function theorem at the zero-discount point gives a unique nearby root. Compactness and the convergence just proved exclude any other equilibria for sufficiently small positive discount. Its Jacobian remains nonsingular. Fix one such positive discount. Compactness, the strict boundary inequalities and nonsingularity preserve these properties under sufficiently small primitive and benefit perturbations, proving the open-set claim. One may fix, for example, $\mu_0$ to the state $(U_1,U_1)$ throughout.

The policy margin is nontrivial. Vary a common unemployment benefit $b$ near one. At state $(U_1,U_1)$, the zero-discount consumption of worker $i$ after its hire is $L_i+w_i+1/2-b/2$, whereas without a hire it is $L_i$. The analytic implicit branch therefore satisfies
\[
 s_i(\beta,b)=0.05\beta\log\frac{L_i+w_i+1/2-b/2}{L_i}
                 +O(\beta^2),\qquad
 \partial_b s_i(\beta,1)=-\frac{0.025\beta}{L_i+w_i}+O(\beta^2)<0
\]
for sufficiently small positive $\beta$. Analyticity makes the remainder differentiable uniformly near $b=1$. Thus benefits change behavior even though their current net receipts cancel when both workers are unemployed.
\end{proof}

The theorem therefore concerns a genuine infinite-horizon funded labor market, including every equilibrium in its declared class. It answers J65-96's local-derivative alternative in that class. Extending the class to unrestricted asset accumulation, behavioral preferences or singular equilibria requires additional analysis.

\section*{Completion Estimate}
100\%
The proof resolves the catalogue's local welfare-derivative alternative
for the declared primitive class. This is the author's complete proof
claim for that target, pending independent expert review.

\begin{thebibliography}{9}
\bibitem{ui} Thomas Le Barbanchon, Johannes Schmieder and Andrea Weber.
\emph{Job Search, Unemployment Insurance, and Active Labor Market Policies}.
RF Berlin Discussion Paper 24/24, October 2024, Sections 3.6.2, 3.7.1 and 3.9.
\url{https://www.rfberlin.com/wp-content/uploads/2024/10/24024.pdf}.
\bibitem{brouwer} L. E. J. Brouwer.
\emph{\"Uber Abbildung von Mannigfaltigkeiten}.
Mathematische Annalen 71 (1911), 97--115.
\url{https://doi.org/10.1007/BF01456931}.
\end{thebibliography}
\end{document}
